\section{Mulțimi, Produse Carteziene și Funcții (04-Oct-1994)}

\subsection*{Teoria Naivă a Mulțimilor}

\begin{definition}
O \textbf{mulțime} este o colecție bine determinată de obiecte distincte, numite elementele mulțimii. Notăm apartenența prin $x \in A$.
\end{definition}

O mulțime poate fi definită prin enumerarea elementelor sale $A = \{a, b, c\}$ sau printr-o proprietate caracteristică:
$$ M = \{x \in A \mid p(x)\} $$
(mulțimea tuturor elementelor din $A$ care satisfac proprietatea $p$).

\begin{definition}[Operații cu mulțimi]\leavevmode
Fie $A$ și $B$ două mulțimi. Definim următoarele operații:
\begin{itemize}
    \item \textbf{Reuniunea:} $A \cup B = \{x \mid x \in A \text{ sau } x \in B\}$
    \item \textbf{Intersecția:} $A \cap B = \{x \mid x \in A \text{ și } x \in B\}$
    \item \textbf{Diferența:} $A \setminus B = \{x \mid x \in A \text{ și } x \notin B\}$
    \item \textbf{Complementara:} Dacă $A \subseteq E$, complementara lui $A$ în raport cu $E$ este $\complement_E A = E \setminus A$.
\end{itemize}
\end{definition}

\begin{remark}
Mulțimea care nu conține niciun element se numește \textbf{mulțimea vidă}, notată $\emptyset$. Mulțimea tuturor submulțimilor lui $A$ se numește \textbf{mulțimea părților} lui $A$, notată $\mathcal{P}(A)$.
\end{remark}

\subsection*{Produsul Cartezian}

\begin{definition}
Fie $A$ și $B$ două mulțimi nevide. Se numește \textbf{pereche ordonată} un ansamblu de două elemente $(a, b)$ unde $a \in A$ și $b \in B$, având proprietatea că $(a, b) = (c, d) \iff a = c$ și $b = d$.
\end{definition}

\begin{definition}
\textbf{Produsul cartezian} a două mulțimi $A$ și $B$ este mulțimea tuturor perechilor ordonate care se pot forma cu elemente din cele două mulțimi:
$$ A \times B = \{(a, b) \mid a \in A \text{ și } b \in B\} $$
\end{definition}

\subsection*{Funcții (Aplicații)}

\begin{definition}
Fie $A$ și $B$ două mulțimi nevide. O \textbf{funcție} (sau aplicație) $f: A \to B$ este o lege de corespondență care asociază fiecărui element $x \in A$ un unic element $y \in B$, notat $y = f(x)$. 
\end{definition}

\begin{definition}[Tipuri de funcții]\leavevmode
O funcție $f: A \to B$ se numește:
\begin{enumerate}
    \item \textbf{Injectivă:} dacă $(\forall) x_1, x_2 \in A, x_1 \neq x_2 \implies f(x_1) \neq f(x_2)$. (Echivalent: $f(x_1) = f(x_2) \implies x_1 = x_2$).
    \item \textbf{Surjectivă:} dacă $(\forall) y \in B, (\exists) x \in A$ astfel încât $f(x) = y$. (Echivalent: $\text{Im} f = B$).
    \item \textbf{Bijectivă:} dacă $f$ este simultan injectivă și surjectivă.
\end{enumerate}
\end{definition}

\begin{definition}
Fie $f: A \to B$ și $g: B \to C$. \textbf{Compunerea} celor două funcții este o nouă funcție $g \circ f: A \to C$ definită prin:
$$ (g \circ f)(x) = g(f(x)), \quad (\forall) x \in A $$
\end{definition}

\begin{proposition}
Compunerea funcțiilor este \textbf{asociativă}. Dacă $f: A \to B$, $g: B \to C$ și $h: C \to D$, atunci:
$$ h \circ (g \circ f) = (h \circ g) \circ f $$
\end{proposition}